\begin{tikzpicture}[unipd/diagram,x=1cm,y=1cm]
  \node at (0,0.65) {Chiamate};
  \node at (6,0.65) {Ritorni};
  \foreach \n/\y/\result in {4/0/24,3/-1/6,2/-2/2,1/-3/1,0/-4/1} {
    \node[unipd/node,minimum width=32mm] (c\n) at (0,\y)
      {\texttt{fattoriale(\n)}};
    \node[unipd/node,minimum width=32mm] (r\n) at (6,\y)
      {\(n=\n\): restituisce \(\result\)};
  }
  \foreach \n/\m in {4/3,3/2,2/1,1/0} {
    \draw[unipd/directed edge] (c\n) -- (c\m);
    \draw[unipd/directed edge] (r\m) -- (r\n);
  }
  \foreach \y/\text in {0/{\(4\cdot6=24\)},-1/{\(3\cdot2=6\)},-2/{\(2\cdot1=2\)},-3/{\(1\cdot1=1\)}} {
    \node at (3,\y) {\text};
  }
  \draw[unipd/directed edge] (c0) -- node[above] {Caso base: \(0!=1\)} (r0);
\end{tikzpicture}
